\chap{Chapter 19. Patching, Vulnerabilities and the Law}
Security is now a legal duty for both the people who make SCADA products and the people who run them. This chapter covers how vulnerabilities are ranked and patched without outages, the European law that binds vendors worldwide, and the Australian law that binds operators. It is a technical summary, not legal advice; Australian facts were checked on 19 November 2026.

\sect{19.1 Ranking vulnerabilities}
\law{19.1}{CVSS severity bands} Vulnerabilities are scored 0 to 10 under the Common Vulnerability Scoring System: 0.1–3.9 low, 4.0–6.9 medium, 7.0–8.9 high, 9.0–10.0 critical.

The score describes the flaw, not your risk. Rank by three questions: is it reachable on our network, is it being exploited, and what could it make the process do?

\wex{19.1} Three findings on one site:
\begin{tbl}{P{0.3\textwidth}P{0.09\textwidth}LP{0.18\textwidth}}
Finding & CVSS & Reachable? Exploited? & Priority\\\midrule
Web server flaw on an HMI behind two firewalls, service unused & 9.8 & Not reachable; no known exploitation & Next outage\\
Remote-access appliance flaw on the internet edge & 7.5 & Reachable; actively exploited & Now\\
Unauthenticated Modbus writes to a PLC & none & Reachable from the control zone; by design & Compensate (Section 17.3)\\
\end{tbl}
The highest score is the lowest priority. Dragos reported that only 2\% of ICS-relevant vulnerabilities in 2025 warranted immediate action.

\law{19.2}{Exposure window} A known vulnerability is open from disclosure until patched, W = t\_patched − t\_disclosed. If exploitation attempts succeed at a constant rate h per day, the chance of compromise is P = 1 − e⁻ʰᵂ. With h = 0.002 per day: 5.8\% if patched within 30 days, 30\% at six months, 52\% at a year. The rate is illustrative; the shape is not. The only variable the site controls is W.

\sect{19.2 Patching without an outage}
Patching used to be avoided for fear that a Windows update would break the SCADA software. Vendors now test operating-system patches against their products and publish the results (Schneider Electric does this for Geo SCADA), which turns patching into a procedure.
\begin{enumerate}
\item Check the vendor's published compatibility for the patch, and the release notes for your version.
\item Take a verified backup of database and configuration.
\item Apply to a test or development system; run through key displays and drivers.
\item Patch the standby server; confirm it resynchronises.
\item Switch over, so the patched server becomes main.
\item Patch the former main; confirm it rejoins as standby.
\item Patch client workstations in groups, never all at once.
\item Record what was installed, where and when, against the work order.
\end{enumerate}
One healthy server is in service throughout (Law 9.6). Measure \emph{patch compliance}, the share of systems in scope at the target level, monthly: what is not measured drifts.

\sect{19.3 The EU Cyber Resilience Act}
The Cyber Resilience Act (Regulation (EU) 2024/2847) makes the manufacturer of a product with digital elements legally responsible for its security throughout a declared support period, generally at least five years. Reporting duties apply from 11 September 2026: early warning of an actively exploited vulnerability within 24 hours, fuller notification within 72 hours, and a final report within 14 days of a fix. Full requirements apply from 11 December 2027, with penalties up to €15 million or 2.5\% of worldwide turnover. Manufacturers must ship secure by default, provide security updates, and keep a software bill of materials (SBOM).

A global vendor does not keep a weaker product for the rest of the world, so Australian sites receive the same hardening and patch cadence. The consequence for operators: once a fix is published, the remaining risk sits with whoever has not installed it.

\sect{19.4 Australian law}
\defn{19.1}{SOCI Act} The Security of Critical Infrastructure Act 2018 covers eleven sectors, including energy, water and sewerage, transport, communications and health. Responsible entities must register assets, report cyber incidents, and run a risk management programme. In water and sewerage it captures larger utilities; the threshold is commonly cited as 100,000 connections (verify for your utility).

\law{19.3}{Incident reporting clock} The clock starts when the entity becomes \emph{aware}, not when the incident began.
\begin{tbl}{P{0.45\textwidth}L}
Impact & Report to ASD's ACSC within\\\midrule
Significant impact on availability & 12 hours\\
Relevant impact (availability, integrity, reliability, confidentiality) & 72 hours\\
Report first made verbally & Written record within 84 hours\\
\end{tbl}
\wex{19.2} At 21:40 on a Friday the on-call technician realises ransomware has stopped the SCADA servers and operators are running the plant by hand. A significant-impact report is due by 09:40 Saturday. Twelve hours overnight leaves no time to work out who is allowed to make the call: the roster, decision maker and reporting steps belong in the incident response plan before the night they are needed. The technician's part is to recognise it quickly and escalate.

\defn{19.2}{CIRMP} The Critical Infrastructure Risk Management Program: cyber and information security, personnel, supply chain, and physical and natural hazards, with a board-approved annual report. For cyber, the entity adopts a recognised framework at a stated level: ISO/IEC 27001, the Essential Eight, the NIST Cybersecurity Framework, C2M2, or the Australian Energy Sector Cyber Security Framework (AESCSF).

\textbf{Cyber Security Act 2024.} Since 30 May 2025, a business with turnover above \$3 million, or the responsible entity for a critical infrastructure asset, must report a ransomware or cyber-extortion payment within 72 hours.

\textbf{ASD's Principles of OT Cyber Security (October 2024).} Safety is paramount; knowledge of the business is crucial; OT data is extremely valuable and needs protecting; segment and segregate OT from all other networks; the supply chain must be secure; people are essential. The third is the one people miss: diagrams, project files and alarm data are what an attacker needs for Stage 2.

\textbf{The Essential Eight in OT.} Application control, restricting administrative privileges, multi-factor authentication for remote access and regular backups (including configuration exports and PLC projects) translate well to SCADA servers. Patching operating systems and applications works for servers in planned windows, rarely for controllers. Office macro settings are mostly irrelevant inside the control zone.

\sect{Exercises}
\begin{enumerate}
\item h = 0.001 per day and W = 90 days. Find the probability of compromise. (8.6\%)
\item 46 of 52 servers are at the target patch level. Find compliance. (88.5\%)
\item An entity becomes aware of a critical incident at 16:30 on a Sunday. When is the report due? (04:30 Monday)
\item List the six ASD principles of OT cyber security. (Section 19.4)
\item Why should a redundant server pair be patched standby first? (One healthy server stays in service throughout)
\end{enumerate}
